%
% claude-sdk-specification-pharo-notes.tex
% Compile: pdflatex claude-sdk-specification-pharo-notes.tex
%

\documentclass[a4paper,10pt,fleqn]{article}
\usepackage[margin=1in]{geometry}
\usepackage[colorlinks=true,linkcolor=blue,citecolor=blue,urlcolor=blue]{hyperref}
\usepackage{booktabs}
\usepackage{array}
\usepackage{tabularx}
\usepackage{longtable}
\usepackage{enumitem}
\usepackage{listings}
\usepackage{xcolor}
\usepackage{amsmath}

\newcolumntype{L}[1]{>{\raggedright\arraybackslash}p{#1}}
\newcolumntype{Y}{>{\raggedright\arraybackslash}X}

\lstset{
  basicstyle=\ttfamily\small,
  keywordstyle=\color{blue!80!black},
  commentstyle=\color{green!50!black},
  stringstyle=\color{red!60!black},
  breaklines=true,
  frame=single,
  framerule=0.3pt,
  rulecolor=\color{gray!50},
  backgroundcolor=\color{gray!5},
  columns=fullflexible,
}

\newcommand{\cls}[1]{\texttt{#1}}
\newcommand{\sel}[1]{\texttt{\##1}}
\newcommand{\pkg}[1]{\texttt{#1}}
\newcommand{\goref}[1]{\textsf{\small [Go: #1]}}

\begin{document}

\title{Claude SDK for Pharo:\\Pharo-Specific Implementation Notes}
\author{Punt Labs}
\date{April 2026 --- v0.3.0-dev}
\hypersetup{
  pdftitle={Claude SDK for Pharo: Pharo-Specific Implementation Notes},
  pdfauthor={Punt Labs},
  pdfsubject={Claude SDK Pharo-Specific Notes},
  pdfcreator={pdflatex}
}
\maketitle

\tableofcontents
\newpage

\subsection*{About this Document}

Pharo-specific notes without an Anthropic counterpart.  Companion
to \emph{claude-sdk-specification.tex}.

\section{Relationship to Go SDK}

The Go SDK (\texttt{anthropic-sdk-go}) is code-generated from Anthropic's
OpenAPI specification via Stainless.  This produces 232 types in
\texttt{message.go} alone, most of which are mechanical variants of a
smaller set of core concepts.

The Smalltalk SDK does not replicate this generated structure.  Instead, it
distills the Go types into a smaller class hierarchy organized by role,
using Smalltalk's polymorphic dispatch and class hierarchy in place of Go's
discriminated unions and struct embedding.

\section{Go to Smalltalk Mapping}

The Go SDK uses flat structs with discriminator fields and \texttt{AsAny()}
type switches for polymorphism.  Smalltalk replaces this with a class
hierarchy and polymorphic message dispatch.

\begin{center}
\begin{tabularx}{\linewidth}{@{}lYY@{}}
\toprule
\textbf{Go Pattern} & \textbf{Smalltalk Equivalent} & \textbf{Example} \\
\midrule
Struct with fields  & Class with instance variables    &
  \cls{ClaudeMessage} \\
Discriminated union (\texttt{AsAny()}) & Abstract class + subclasses &
  \cls{ClaudeContentBlock} hierarchy \\
\texttt{param.Opt[T]} (optional) & \texttt{nil} or value &
  Standard Smalltalk \\
\texttt{json.RawMessage} & Raw JSON string &
  String or Dictionary \\
Param vs Response types & Single class (Smalltalk is dynamic) &
  \cls{ClaudeMessage} for both \\
Constants (\texttt{type X string}) & Class-side methods returning symbols &
  \cls{ClaudeModel class>>sonnet} \\
\bottomrule
\end{tabularx}
\end{center}

\section{JSON Key Mapping}

Wire keys are snake\_case strings (\texttt{max\_tokens}).  Smalltalk instance
variables are camelCase (\texttt{maxTokens}).  Every type class defines a
class-side method \sel{jsonKeyMap} returning a Dictionary mapping instance
variable names (Symbols) to wire key strings:

\begin{lstlisting}[language={},title={JSON key mapping}]
ClaudeMessage class>>jsonKeyMap
    ^ {
        #stopReason -> 'stop_reason'.
        #stopSequence -> 'stop_sequence'.
        #inputTokens -> 'input_tokens'.
    } asDictionary
\end{lstlisting}

The \sel{asJson} method iterates instance variables, looks up the wire key
from the map (defaulting to the variable name if absent), and omits nil
values.  The class-side \sel{fromJson:} reverses the map to populate
instance variables from the parsed Dictionary.

\section{Helper Constructors}

Class-side factory methods on content block classes provide concise construction:

\begin{lstlisting}[language={},title={Helper constructors}]
ClaudeTextBlock text: 'Hello, world'.
ClaudeImageBlock base64: encodedData mediaType: 'image/png'.
ClaudeToolUseBlock id: 'tu_123' name: 'get_weather' input: inputDict.
ClaudeToolResultBlock toolUseId: 'tu_123' content: 'Sunny, 22C'.
ClaudeToolResultBlock toolUseId: 'tu_123' content: 'Not found' isError: true.
ClaudeThinkingBlock thinking: 'Let me reason...' signature: 'sig...'.
ClaudeRedactedThinkingBlock data: 'redacted-data-string'.
ClaudeDocumentBlock fromBase64: 'JVBERi0xLjQ='.
ClaudeDocumentBlock fromFileId: 'file-abc123'.
ClaudeDocumentBlock fromText: 'Plain text document content'.
ClaudeDocumentBlock fromUrl: 'https://example.com/doc.pdf'.
ClaudeImageBlock fromUrl: 'https://example.com/image.png'.
ClaudeImageBlock fromFileId: 'file-img456'.
\end{lstlisting}

\section{JSON Serialization}

The Go SDK uses a custom JSON codec (\texttt{internal/apijson/}) with
reflection-based encoder/decoder caches.  Smalltalk does not need this
--- STONJSON handles serialization directly.

\subsection{Conventions}

\begin{itemize}
  \item \textbf{Request serialization:} Each type implements \sel{asJson}
    returning a Dictionary.  \texttt{STONJSON toString:} converts to
    wire format.
  \item \textbf{Response deserialization:} \texttt{STONJSON fromString:}
    returns a Dictionary (specifically \texttt{OrderedDictionary} in
    Pharo~12).  Each type's class-side \sel{fromJson:} constructs
    typed objects.
  \item \textbf{Nil omission:} \sel{asJson} omits nil-valued fields.
    The Go SDK uses \texttt{param.Opt[T]} for this; Smalltalk uses nil
    naturally.
  \item \textbf{Key format:} Wire keys are snake\_case strings
    (\texttt{max\_tokens}).  Smalltalk instance variables are camelCase
    (\texttt{maxTokens}).  The \sel{jsonKeyMap} class-side method defines
    the mapping (see the JSON Key Mapping section above).
  \item \textbf{Unknown fields:} When deserializing, unknown JSON keys are
    silently ignored.  This provides forward compatibility when the API adds
    new fields.
\end{itemize}

\section{Pharo Threading Model (Streaming)}

Pharo uses cooperative scheduling within a priority level.  \texttt{[ ... ]
fork} spawns a \texttt{Process} (green thread) that runs cooperatively with
other processes at the same priority.  Processes at the same priority do
\emph{not} run in parallel --- they yield at I/O points, explicit yields,
and semaphore waits.  Processes at different priorities are preemptive (higher
priority preempts lower).

\textbf{Rule:} streaming network I/O must run in a forked Process.

\begin{lstlisting}[language={},title={Simple streaming pattern (buffered)}]
| events sem |
events := OrderedCollection new.
sem := Semaphore new.

[ [ client
    streamMessage: params
    do: [ :event | events add: event ].
  sem signal ]
    on: ClaudeApiError
    do: [ :e | events add: e. sem signal ] ]
  forkAt: Processor userBackgroundPriority
  named: 'Claude stream'.

"Consumer waits for completion, then processes all events:"
sem wait.
events do: [ :each |
    each isKindOf: ClaudeApiError
        ifTrue: [ self handleError: each ]
        ifFalse: [ self handleEvent: each ] ]
\end{lstlisting}

The SDK provides \sel{streamMessage:do:} as a synchronous method.  The
consumer is responsible for forking.  The SDK documents this requirement
but does not enforce it --- forcing a fork inside the SDK would prevent
consumers from choosing their own concurrency strategy.

For incremental delivery, use the real-time mode pattern (SharedQueue +
forked socket reader) described in the main specification's Streaming
Messages section.

\section{Porting from Go SDK}

The reference implementation is Anthropic's Go SDK
(\texttt{anthropic-sdk-go}), code-generated from the OpenAPI specification
via Stainless.  The Go SDK produces 232 types in \texttt{message.go}
alone --- most are mechanical variants of a smaller set of core concepts
(param vs.\ response structs, union wrappers, JSON marshalers).

The porting strategy is \emph{distillation, not transliteration}.
Smalltalk's class hierarchy and polymorphic dispatch replace Go's flat
structs and discriminated unions.  A single \cls{ClaudeContentBlock}
hierarchy with 16 subclasses replaces 30+ Go \texttt{ContentBlock*}
types.  \cls{ClaudeMessage} serves as both request parameter and
response --- Go needs separate \texttt{Message} and
\texttt{MessageParam} structs because the type system cannot share them.

Each Go type is mapped once via the type inventory (see the Go SDK Type
Inventory section below).  When a Go type group collapses to a single
Smalltalk class, the mapping is recorded and the redundant Go variants
are not ported.  This keeps the Smalltalk class count proportional to
the number of \emph{concepts}, not the number of generated types.

The reference SDK is kept at \texttt{.tmp/anthropic-sdk-go/} (gitignored)
and consulted for every new API feature to ensure wire-format
compatibility.

\section{Postern Bootstrap}

All development happens inside a live Pharo~12 image via Postern, an
HTTP eval server running on \texttt{localhost:\$\{EVAL\_PORT:-8422\}}.
The eval server accepts Smalltalk expressions as plain-text POST bodies
and returns printed results.  This provides interactive class definition,
method compilation, test execution, lint checks, and full introspection
--- all driven from Claude Code's Bash tool or any HTTP client.

The image is disposable.  \texttt{make rebuild} creates a fresh image
from scratch, loading all source from Tonel files in \texttt{src/}.
\texttt{make stop} kills without saving --- saving would capture socket
state that causes errors on restore.  The source of truth is always
Tonel on disk, committed via Iceberg (libgit2) from inside the image.

Key Postern targets: \texttt{make setup} (download Pharo~12 + load
packages), \texttt{make start} (launch with eval server),
\texttt{make test} (run SDK test suite), \texttt{make lint} (Renraku
lint on all classes), \texttt{make drift} (compare in-image methods
vs.\ on-disk Tonel), \texttt{make filein} (reload Tonel into running
image).

Multiple images can run simultaneously on different ports via the
\texttt{EVAL\_PORT} environment variable, each with its own PID file
and log file.

\section{Agent-Assisted Development}

Development uses ethos mission pipelines to delegate work to specialist
agents.  The standard pipeline has five stages: design, implement, test,
review, document.  Each stage produces an artifact that feeds the next
via \texttt{inputs\_from} dependency wiring.

Implementation work is dispatched to a Smalltalk specialist agent (kwb)
scoped by a mission contract that specifies the write-set (which Tonel
files may be touched), success criteria (\texttt{make lint} clean,
specific test classes passing, Iceberg ref sync), round budget, and
evaluator.  Every contract is peer-reviewed by an architect agent before
the worker sees it.

The SDK's own \cls{ClaudeExchange} agent runtime is used during
development --- the implementing agent runs inside the Pharo image via
the eval server, compiling methods and running tests through the same
SDK classes it is building.  This in-image dogfooding exercises
interaction patterns that no external test harness can simulate:
streaming under load, tool dispatch cycles, error recovery across
conversation turns.

Quality gates are layered: the worker runs per-class tests via the eval
server, the leader runs \texttt{make lint} from CLI (which catches
class-level Renraku rules that per-method probes miss), and local review
agents scan the full branch diff before any PR is created.

\section{Package Health}

Package structure is assessed using the Ducasse/Abdeen modularization
metrics from WCRE'11 (\emph{Package Fingerprints}).  The key metric is
IIPU (Index of Inter-Package Usage) --- the ratio of intra-package
references to total references.  Higher IIPU means better encapsulation.

The SDK tracks IIPU across restructuring milestones.  The original flat
\texttt{Claude-SDK-*} layout scored IIPU 0.458 (heavy cross-package
coupling).  After the Messaging/Agent layer split and the introduction
of \texttt{Claude-Messaging-Tools} and \texttt{Claude-Messaging-Files}
as separate packages, IIPU rose to 0.838.  The \texttt{/package-health}
assessment skill provides repeatable analysis.

Naming discipline follows Kent Beck's Intention Revealing Selector
pattern, scaled from methods to packages.  Package names follow the
Pharo idiom \texttt{Family-SubFamily-Module} (e.g.,
\texttt{Claude-Messaging-Types}, \texttt{Claude-Agent-Tools-Pharo}).
The layered names make the dependency hierarchy explicit: everything in
\texttt{Claude-Agent-*} depends on \texttt{Claude-Messaging-*}, never
the reverse.  Tool names use PascalCase compound names matching Claude
Code's convention (\texttt{BrowseClass}, \texttt{CompileMethod},
\texttt{RunTests}).

Every package name change and every new package is reviewed by an
architect agent for consistency with the naming scheme before
implementation begins.

\section{Package Architecture}
\label{notes:package-architecture}

\subsection{Package Inventory}

\begin{center}
\begin{tabularx}{\linewidth}{@{}lrY@{}}
\toprule
\textbf{Package} & \textbf{Classes} & \textbf{Responsibility} \\
\midrule
\pkg{Claude-Messaging-Types}      & 48 & \cls{ClaudeMessage} and decomposition
  (turn, stop details, request context, container), content block hierarchy
  (16 subclasses including \cls{ClaudeDocumentBlock} with source hierarchy),
  tool types, model, usage, cache control, thinking params, sampling params,
  request options, output format; JSON serialization; helper constructors \\
\pkg{Claude-Messaging-Errors}    & 10 & \cls{ClaudeApiError} hierarchy (8 concrete
  subclasses keyed on HTTP status), error factory, \cls{ClaudeInvalidRequestError} \\
\pkg{Claude-Messaging-Streaming} &  2 & \cls{ClaudeSSEDecoder} (buffered and real-time
  SSE parsing), \cls{ClaudeStreamingSocket} (raw TLS with chunked
  transfer-encoding) \\
\pkg{Claude-Messaging-Client}    &  1 & \cls{ClaudeClient} --- creation, configuration,
  HTTP request lifecycle, options, retry logic, Files API \\
\pkg{Claude-Messaging-Tools}     & 10 & \cls{ClaudeServerTool} hierarchy (8 concrete
  server-side tools), name-keyed \cls{TypeRegistry} \\
\pkg{Claude-Messaging-Files}     &  4 & \cls{ClaudeFileMetadata},
  \cls{ClaudeFilePage}, \cls{ClaudeFileListParams},
  \cls{ClaudeDeletedFile} \\
\pkg{Claude-Messaging-Examples}  &  7 & Example scripts for messaging API usage \\
\bottomrule
\end{tabularx}
\end{center}

\subsection{Dependency Flow}

Dependencies flow in strict layers:

\begin{verbatim}
  Messaging-Types  Messaging-Errors  Messaging-Streaming
        |               |                |
  Messaging-Tools  Messaging-Files       |
        |               |                |
        +-------+-------+-------+--------+
                |
          Messaging-Client
\end{verbatim}

\noindent Messaging-Client depends on Types, Errors, Streaming, Tools
(server-side tool types), and Files (file metadata types).

The Agent layer (\pkg{Claude-Agent-*}) depends on Messaging-Client and
is specified in the \emph{Claude Agent SDK for Pharo --- Technical
Specification}.

\section{Concurrency Model}

\subsection{The Go SDK is Synchronous}

The Go SDK has no goroutines in its public API.  All methods are blocking.
\texttt{http.Client} manages connection pooling internally.  The Smalltalk
SDK mirrors this: all public methods are synchronous and blocking.

\subsection{What Runs Where}

\begin{center}
\begin{tabularx}{\linewidth}{@{}lYl@{}}
\toprule
\textbf{Operation} & \textbf{Blocking?} & \textbf{Safe for UI?} \\
\midrule
Create client              & No  & Yes \\
Build message params       & No  & Yes \\
\sel{sendMessage:}         & Yes (network I/O) & \textbf{No --- fork} \\
\sel{streamMessage:do:}    & Yes (network I/O) & \textbf{No --- fork} \\
\sel{listModels}           & Yes (network I/O) & \textbf{No --- fork} \\
\sel{getModel:}            & Yes (network I/O) & \textbf{No --- fork} \\
\sel{countTokens:}         & Yes (network I/O) & \textbf{No --- fork} \\
Parse JSON response        & No  & Yes \\
Tool handler execution     & Depends & Fork if I/O-bound \\
Retry backoff sleep        & Yes & \textbf{No --- fork} \\
\bottomrule
\end{tabularx}
\end{center}

\subsection{Safe Patterns}

\begin{description}[style=nextline]
  \item[Fork + Semaphore]
    Fork a Process, signal a Semaphore on completion.  Consumer waits on
    the Semaphore.  Use when the consumer needs the full result before
    proceeding.

  \item[Fork + SharedQueue]
    Fork a Process, enqueue results into a SharedQueue.  Consumer drains
    the queue incrementally.  Use for streaming events where partial results
    are useful.  The consumer reads from the queue in a loop, not after
    a semaphore wait.

  \item[Fork + callback]
    Fork a Process, call a block on each event.  The block must not touch
    Morphic directly --- use \texttt{WorldState defer:} to schedule UI
    updates on the UI thread.
\end{description}

\subsection{Unsafe Patterns}

\begin{itemize}
  \item Calling \sel{sendMessage:} directly from a Morphic event handler.
  \item Calling \sel{streamMessage:do:} without forking.
  \item Modifying Morphic submorphs from a background Process without
    \texttt{WorldState defer:}.
  \item Waiting on a Semaphore from a SharedQueue-based streaming pattern
    (blocks until complete, defeating incremental delivery).
\end{itemize}

\section{Smalltalk Idioms}

Kent Beck's \emph{Smalltalk Best Practice Patterns} guide the SDK design.

\subsection{Composed Method (Pattern 1)}

Every method does one identifiable thing.  Large methods are decomposed
into helper messages.  The Go SDK's \texttt{ExecuteNewRequest} (80+
lines) becomes several composed methods in Smalltalk:
\sel{buildRequest:}, \sel{executeWithRetry:}, \sel{parseResponse:}.

\subsection{Intention Revealing Selector (Pattern 2)}

Method names describe \emph{what}, not \emph{how}:

\begin{itemize}
  \item \sel{sendMessage:} not \sel{postToMessagesEndpoint:}
  \item \sel{textContent} not \sel{collectTextBlocksAndJoin}
  \item \sel{toolUseBlocks} not \sel{selectContentBlocksOfTypeToolUse}
\end{itemize}

\subsection{Pluggable Behavior (Pattern 13)}

Use blocks for extensibility:

\begin{itemize}
  \item \sel{streamMessage:do:} takes a callback block.
  \item Tool runner handlers are blocks.
  \item Retry policy could be a pluggable strategy block.
\end{itemize}

\subsection{Collecting Parameter (Pattern 14)}

Build results incrementally.  The message accumulator is a collecting
parameter: it receives streaming events one at a time and constructs the
full message progressively.

\subsection{Method Object (Pattern 15)}

When a method gets too complex, extract it into a separate class whose
instances represent a single execution of the algorithm.  In this SDK,
the following are already designed as method objects (separate classes,
not methods buried in \cls{ClaudeClient}):

\begin{itemize}
  \item \cls{ClaudeSSEDecoder} --- encapsulates SSE parsing state
  \item \cls{ClaudeExchange} --- encapsulates the tool dispatch loop
    with iteration count and state inspection (see the
    \emph{Claude Agent SDK Specification})
\end{itemize}

The retry loop is part of \cls{ClaudeClient}'s request lifecycle.  If
retry logic grows more complex (e.g., pluggable retry policies), it
should be extracted into a \cls{ClaudeRetryPolicy} method object.

\section{Go SDK Type Inventory (Condensed)}
\label{notes:type-inventory}

The Go SDK's \texttt{message.go} contains 232 types.  Most are generated
variants.  The following table maps every Go type group to the Smalltalk
class that replaces it.

\begin{center}
\small
\begin{longtable}{@{}llp{4.5cm}@{}}
\toprule
\textbf{Go Type Group} & \textbf{Smalltalk Class} & \textbf{Notes} \\
\midrule
\endfirsthead
\toprule
\textbf{Go Type Group} & \textbf{Smalltalk Class} & \textbf{Notes} \\
\midrule
\endhead
Message, MessageJSON       & \cls{ClaudeMessage}        & Response + param;
  includes \texttt{type} field \\
MessageParam               & via \sel{asParam}          & Converter method \\
MessageNewParams           & \cls{ClaudeMessageRequest}  & Request builder \\
ContentBlockUnion (12+ variants) & \cls{ClaudeContentBlock} subclasses &
  Polymorphic hierarchy \\
ContentBlockParamUnion     & Same classes, \sel{asJson}  & Shared \\
TextBlock (with Citations) & \cls{ClaudeTextBlock}      & \texttt{text} +
  \texttt{citations} \\
ThinkingBlock              & \cls{ClaudeThinkingBlock}  & \texttt{thinking},
  \texttt{signature} \\
RedactedThinkingBlock      & \cls{ClaudeRedactedThinkingBlock} & \texttt{data} \\
ServerToolUse, WebSearch*, etc. & Server result subclasses & Read-only,
  raw JSON stored \\
Usage, MessageDeltaUsage   & \cls{ClaudeUsage}          & Single class \\
ToolParam, ToolUnionParam  & \cls{ClaudeTool}           & Unified \\
ToolInputSchemaParam       & \cls{ClaudeToolInputSchema} & Schema builder \\
ToolResultBlockParam       & \cls{ClaudeToolResultBlock} & In hierarchy \\
ToolChoiceUnionParam       & Dictionary                 & Dynamic structure \\
ThinkingConfigParamUnion   & \cls{ClaudeThinkingParams} & Typed class; 3 modes:
  enabled, adaptive, disabled \\
StopReason (enum)          & Symbol                     & 6 values \\
Model (enum)               & \cls{ClaudeModel} class methods & Factory
  pattern; alias + dated \\
ModelInfo                  & \cls{ClaudeModelInfo}      & From Models API \\
MessageTokensCount         & \cls{ClaudeMessageTokensCount} & From CountTokens \\
MessageStreamEventUnion    & \cls{ClaudeStreamEvent} subclasses & Polymorphic \\
*Delta types (5)           & Embedded in event subclasses & Simplified \\
RawContentBlockDeltaUnion  & Discriminated in accumulator & 5 delta subtypes \\
Base64PDFSource            & \cls{ClaudeDocumentBlock}       & Document content block
  with \cls{ClaudeDocumentSource} hierarchy (5 source types) \\
CacheControlEphemeralParam & \cls{ClaudeCacheControl}   & Typed class with optional TTL \\
MetadataParam              & Dictionary                 & Dynamic structure \\
ServiceTier                & String                     & \texttt{'standard'} etc. \\
MessageBatch*              & Deferred                   & v2 \\
Beta*                      & Deferred                   & v2 \\
Bedrock*, Vertex*          & Out of scope               & --- \\
Completion*                & Out of scope               & Legacy \\
\bottomrule
\end{longtable}
\normalsize
\end{center}

\section{Pharo Threading Quick Reference}

Pharo uses cooperative scheduling.  Processes at the same priority level
run one at a time, yielding at I/O waits, explicit yields (\texttt{Processor
yield}), and semaphore waits.  Higher-priority processes preempt lower ones.
This is \emph{not} parallel execution --- a CPU-bound process at user
priority will starve other user-priority processes until it yields.

\begin{center}
\begin{tabularx}{\linewidth}{@{}lY@{}}
\toprule
\textbf{Primitive} & \textbf{Usage} \\
\midrule
\texttt{[ ... ] fork}  & Spawn a Process at the same priority as the
  caller.  Does not inherit the caller's stack. \\
\texttt{[ ... ] forkAt:named:} & Spawn at a specific priority with a name.
  Use \texttt{Processor userBackgroundPriority} for SDK work. \\
\texttt{Semaphore new}  & Binary signal/wait.  Use for ``notify when
  done.'' \\
\texttt{SharedQueue new} & Thread-safe FIFO.  Use for producer/consumer
  event passing.  \texttt{next} blocks until an item is available. \\
\texttt{Mutex new}      & Mutual exclusion.  Use for shared mutable
  state (rare in this SDK). \\
\texttt{Delay forSeconds:} & Non-busy wait.  Yields to other processes
  during the delay.  Used for retry backoff. \\
\texttt{WorldState defer: [ ... ]} & Schedule a block to run on the
  Morphic UI thread.  Required for any Morph modification from a
  background Process. \\
\texttt{Process>>terminate} & Kill a running process.  Use for
  cancellation (e.g., abort a streaming request).  Closes any
  associated sockets. \\
\bottomrule
\end{tabularx}
\end{center}

\end{document}
